\xiaosi

NCP的布线拓扑由6个离散参数决定（见2.2节），参数间存在不等式约束。为了自动发现优于手工设计的布线方案，本文设计了一套基于进化策略的搜索框架。

\subsubsection{搜索空间与约束}

搜索空间定义为6维整数向量$g = (n_i, n_c, f_s, f_i, r_c, f_m)$，各分量的取值范围为：$n_i \in [2,20]$，$n_c \in [2,16]$，$f_s \in [1,8]$，$f_i \in [1,8]$，$r_c \in [0,8]$，$f_m \in [1,8]$。约束条件为$f_s \leq n_i$，$f_i \leq n_c$，$f_m \leq n_c$。违反约束的基因组在每次遗传操作后通过裁剪修复。

\subsubsection{安全感知适应度函数}

与仅优化奖励的传统方法不同，本文设计了安全感知的多目标适应度函数：

\begin{equation}
    F(g) = \bar{R}(g) - \alpha \cdot \bar{C}(g)
    \label{eq:fitness}
\end{equation}

其中$\bar{R}(g)$和$\bar{C}(g)$分别为候选布线$g$在训练后的平均累计奖励和平均碰撞率，$\alpha = 10$为安全惩罚系数。该设计使搜索不仅追求高奖励，还主动降低碰撞率，直接回应自动驾驶的安全性需求。

\subsubsection{搜索流程}

具体搜索参数为：种群大小16，迭代15代，锦标赛大小3，精英保留2个，交叉率0.5，变异率0.3。每个候选个体在highway-v0环境上快速训练8,000步后评估适应度。搜索过程中，一代内的所有个体在多GPU上并行评估，每代耗时约30分钟。

\subsubsection{搜索结果}

经过15代240次评估，搜索发现的最优布线为$g^* = (3, 7, 3, 5, 0, 2)$，即仅需3个中间神经元、7个命令神经元（共15个神经元，含5个运动神经元），远少于手工设计的25个。一个值得注意的发现是$r_c = 0$，即最优拓扑\textbf{不需要命令层内的循环连接}。这表明在highway-v0的决策时间尺度上，LTC的ODE动力学本身已提供了足够的短时记忆能力，额外的循环连接反而增加了不必要的复杂度。

搜索出的拓扑在highway-v0上取得了与手工设计相近的奖励（33.88对34.94），但在未参与搜索的merge-v0和roundabout-v0上性能有所下降（11.12对15.00和5.47对6.63），说明单环境搜索的拓扑存在过拟合风险，未来可通过多环境联合搜索改善泛化能力。
